\documentclass[11pt,letterpaper,sans]{moderncv}

% Geometry
\usepackage[scale=0.9]{geometry}

% ModernCV style
\moderncvstyle{banking}

% Fonts
\usepackage{lmodern}

\usepackage{fontawesome5}

% Colors
\usepackage[dvipsnames]{xcolor}
\definecolor{myred}{RGB}{200,0,0}  % Bright red

% --------------------------
% Red section headings
% --------------------------
\makeatletter
\renewcommand*{\section}[1]{%
  \vspace{1.5ex}%
  {\color{myred}\bfseries\Large #1}\par%
  \vspace{0.5ex}%
}
\makeatother

% --------------------------
% Personal data
% --------------------------
\name{\textcolor{myred}{Mohamed}}{\textcolor{myred}{Abuella}}
\address{Newcastle}{United Kingdom}
\phone[mobile]{+44~(745) (955)~6542}
\email{mhdabuella@gmail.com}
\homepage{mohamedabuella.github.io}
%\homepage{mohamedabuella.github.io | \href{https://mohamedabuella.github.io/cv/}{\textcolor{blue}{Detailed CV Page}}}
%\social[Skype][join.skype.com/invite/REqE79GkflGF]{Skype: mohammed\_abuella}
\social[linkedin]{mohamed-abuella}
\social[twitter]{mhdabuella}
\social[github]{MohamedAbuella}
%\social[github]{mhdella}
\social[GS][scholar.google.com/citations?user=nJbFUkcAAAAJ&hl=en]{Scholar: Mohamed Abuella}

\begin{document}

%--------------------------
% Make title
%--------------------------
\makecvtitle
\vspace{-2ex}

%--------------------------
\section{Summary}
%I am interested in modernizing the grid and optimizing its integration of renewables by applying descriptive, predictive, and prescriptive analytics.
%My broader interest encompasses the utilization of Artificial Intelligence to foster Sustainability. Looking for opportunities to transfer and improve my knowledge and skills.
%The technical skills are presented as Acquired Expertise in the Experience section.

%My research develops data-driven decision analytics for integrated energy systems to support the transition toward net-zero energy infrastructures. I combine artificial intelligence, optimization, and game-theoretic modeling to analyze the planning and operation of modern energy systems with high penetration of renewable resources and emerging technologies such as hydrogen.

%My work spans energy forecasting, spatiotemporal analytics, and strategic infrastructure planning, with applications in renewable integration, energy efficiency, and sector coupling. Through interdisciplinary research that bridges power systems engineering, artificial intelligence, and policy-oriented modeling, I aim to develop analytical tools that support resilient, efficient, and sustainable energy transitions.

%I am particularly interested in advancing AI-enabled decision support for complex socio-technical energy systems and collaborating across academia, industry, and policy stakeholders to accelerate the deployment of sustainable energy solutions.


%My research develops data-driven decision analytics for integrated energy systems to support the transition to net-zero infrastructures. I combine optimization and game-theoretic modeling to analyze the planning and operation of modern energy systems with high penetration of renewable resources and emerging technologies such as maritime and hydrogen.
%My work spans energy forecasting, spatiotemporal analytics, and strategic infrastructure planning, with applications in renewable integration, energy efficiency, and sector coupling.
%Looking for opportunities to transfer and improve my knowledge and skills. \\


Electrical engineer by training, traditionally interested in mathematical and computational analysis, modeling, and optimization, and data analytics for integrated energy systems planning and operations.

As a researcher, I focus on modernizing the electric grid and optimizing the integration of distributed energy resources by applying descriptive, predictive, and prescriptive analytics to improve energy efficiency and support the transition to net-zero infrastructures. My work includes recent applications in hydrogen and maritime shipping, with a broader interest in utilizing Artificial Intelligence to foster sustainability. Seeking to transfer and expand my expertise to bridge the gap between advanced research and real-world application.

\textbf{Keywords:} \textit{Energy systems modeling and analytics, planning and operation, renewable energy and hydrogen integration, optimization and game theory, energy transition policy, net-zero pathways, AI/ML for energy forecasting and efficiency.}

%\textbf{Keywords:} \textit{Power systems, optimization and modeling, AI/ML for energy systems, energy policy and game theory, hydrogen and integrated energy systems, energy systems architecture and strategy.}

%with a growing focus on energy systems architecture and strategy.%

%{\color{myred}$\rightarrow$} For more details, see my online CV and professional profile:
{\faLink}\, For more details, see my online CV and professional profile in
\href{https://mohamedabuella.github.io/cv/}{\textcolor{blue}{mohamedabuella.github.io/cv}}.

%--------------------------
\section{Education}
\cventry{2014--2018}{Ph.D in Electrical Engineering}{University of North Carolina at Charlotte (UNCC)}{USA}{GPA 4.0}{}
\cventry{2010--2012}{M.Sc in Electrical and Computer Engineering}{Southern Illinois University at Carbondale (SIUC)}{USA}{GPA 4.0}{}
\cventry{2002--2008}{DipHE in Instrumentation {82\%}, and B.Tech Electromechanical Engineering}{Higher Polytechnic Institute \& College of Industrial Technology at Misurata}{Libya}{86\%}{}

%--------------------------
\section{Experience}
\cventry{May 2024--May 2026}{Northumbria University, Department of Mathematics, Physics and Engineering}{Research Fellow}{UK}{}{
Supporting the delivery of the “Hydrogen Integration for Accelerated Energy Transitions (HI-ACT)” research project.
\newline{}
See more details in \href{https://mohamedabuella.github.io/blog/2026/05/07/HI_ACT_Project_Hydrogen_Integration_for_Accelerated_Energy_Transition_in_GB}
{\textcolor{blue}{HI-ACT Project}}.
\newline{}
\cvlistitem{Acquired Expertise: \emph{Integrated Energy Systems, Game Theory, Strategic \& Policy Modeling, Hydrogen Integration, Energy Transition}}}

\cventry{April 2022--April 2024}{Halmstad University}{Researcher}{Sweden}{}{
Postdoctoral Researcher at the Center for Applied Intelligent Systems Research (CAISR).
Research on AI for Sustainability applying Machine Learning techniques.
\newline{}
For more details in
\href{https://mohamedabuella.github.io/blog/2024/03/24/iHelm_Project_for_Improving_Energy_Efficiency_in_Short_Sea_Shipping}
{\textcolor{blue}{iHelm Project}}.
\newline{}
\cvlistitem{Acquired Expertise: \emph{Integrated Industry-Academia Collaboration, Big-Data \& Spatiotemporal Analysis, Explainable AI (XAI), Research Methodologies \& Project Supervision.}}}

\cventry{February 2020--March 2022}{College of Industrial Technology at Misurata}{Lecturer}{Libya}{}{
Taught Electrical Circuits, Electrical Measurements, Math 101.
\newline{}
\cvlistitem{Acquired Expertise: \emph{Curriculum Revision \& Preparing, Dedication, Listening, "Try to Modeling the Student's Way of Thinking". }}}

\cventry{2014--2019}{Energy Production and Infrastructure Center (EPIC) at UNC Charlotte}{Research Assistant}{USA}{}{
Statistical and Predictive Analytics to Modernize the Grid and Optimize its Integration of Renewables, focusing on Solar Energy Resources.
\newline{}
For more details in
\href{https://mohamedabuella.github.io/blog/2021/06/30/MSc_PhD_Research_in_Energy_Analytics} {\textcolor{blue}{Graduate Research Portfolio}}.
\newline{}
\cvlistitem{Acquired Expertise: \emph{Energy Analytics, Energy Markets, Renewable Energy \& Supply Chain, Machine Learning.}}}

\cventry{2010--2012}{Department of Electrical and Computer Engineering at SIUC}{M.Sc Research}{USA}{}{
Optimization for Electric Power Systems Including Wind Power.
\newline{}
For more details in
\href{https://mohamedabuella.github.io/blog/2021/06/30/MSc_PhD_Research_in_Energy_Analytics} {\textcolor{blue}{Graduate Research Portfolio}}.
\newline{}
\cvlistitem{Acquired Expertise: \emph{Power Systems Analysis, Operation and Planning, Systems Optimization, Smart Grid.}}}

\cventry{2008--2009}{College of Industrial Technology at Misurata}{Teaching Assistant and Lab Instructor}{Libya}{}{
Taught Mathematics, Power Systems Analysis, and PLC.
\newline{}
\cvlistitem{Acquired Expertise: \emph{Teaching, Tutorials, Lab Modeling \& Simulations.}}}

\cventry{2000--2008}{Residential Electrical Wiring and Water \& Wastewater Company}{Electrical Technician}{Libya}{}{
Wiring and maintain electrical control equipment.
\newline{}
\cvlistitem{Acquired Expertise: \emph{Electrical Wiring \& Installations, Maintenance \& Operation.}}}

%--------------------------
\section{Recognitions}
\cvitemwithcomment{Best Paper Award (Third Prize)}{RPG 2025 – 14th Int. Conf. on Renewable Power Generation, Shanghai, China}{2025}
\cvitemwithcomment{Outstanding Reviewer}{IEEE Transactions on Sustainable Energy}{2017}
\cvitemwithcomment{Third Prize for Student Papers}{The 47th North American Power Symposium}{2015}
\cvitemwithcomment{The Institute Prize for team of Energy Analytics Lab at UNCC}{Global Energy Forecasting Competition}{2014}
\cvitemwithcomment{The 1st Place}{Department of Electromechanical Engineering at College of Industrial Technology}{2008}

%--------------------------

\section{Publications}
Authored numerous peer-reviewed publications (all as first author). For the complete list, see my Google Scholar profile:
\href{https://scholar.google.com/citations?user=nJbFUkcAAAAJ&hl=en}{\underline{\textcolor[rgb]{0.00,0.07,1.00}{Mohamed Abuella}}.}
Accumulating over 700 citations and current H-index is \textbf{9}.
my \href{https://orcid.org/0000-0001-8018-9998}{\underline{\textcolor[rgb]{0.00,0.07,1.00}{ORCID}}.}

%--------------------------
\subsection*{Journal Articles}

1. M. Abuella, A. Allahham, S. L. Walker,
"Game Theory Approaches to Hydrogen Infrastructure Investment Planning in Great Britain: A Comparative Analysis of Competitive and Cooperative Frameworks,"
\textit{International Journal of Hydrogen Energy}, vol.~220, p.~154064, 2026.
DOI: \href{https://doi.org/10.1016/j.ijhydene.2026.154064}{\textcolor{blue}{10.1016/j.ijhydene.2026.154064}}.

2. M. Abuella, A. Allahham, S. L. Walker,
"A Cooperative Planning Framework for Hydrogen Blending in Great Britain's Integrated Energy System,"
\textit{Energies}, vol.~19, no.~9, p.~2018, 2026.
DOI: \href{https://doi.org/10.3390/en19092018}{\textcolor{blue}{10.3390/en19092018}}.

3. M. Abuella, A. Allahham, N. A. Rufa'I, S. L. Walker,
"Hydrogen-Based Decarbonisation Strategies for Residential Heating: An Energy Efficiency and Conservation Analysis in the North of Tyne Region,"
\textit{Energies}, vol.~18, no.~23, p.~6237, 2025.
DOI: \href{https://doi.org/10.3390/en18236237}{\textcolor{blue}{10.3390/en18236237}}.

4. M. Abuella, H. Fanaee, S. Nowaczyk, S. Johansson, E. Faghani,
"Time-Series Analysis Approach for Improving Energy Efficiency of Fixed-Route Passenger Vessel in Short-Sea Shipping,"
\textit{Ocean Engineering}, vol.~334, p.~121555, 2025.
DOI: \href{https://doi.org/10.1016/j.oceaneng.2025.121555}{\textcolor{blue}{10.1016/j.oceaneng.2025.121555}}.

5. M. Abuella, M. Atoui, S. Nowaczyk, S. Johansson, E. Faghani,
"Spatial Clustering Approach for Vessel Path Identification,"
\textit{IEEE Access}, vol.~12, pp.~66248--66258, 2024.
DOI: \href{https://doi.org/10.1109/ACCESS.2024.3399116}{\textcolor{blue}{10.1109/ACCESS.2024.3399116}}.

6. M. Abuella and B. Chowdhury,
"Forecasting of Solar Power Ramp Events: A Post-Processing Approach,"
\textit{Renewable Energy}, vol.~133, pp.~1380--1392, 2019.
DOI: \href{https://doi.org/10.1016/j.renene.2018.09.005}{\textcolor{blue}{10.1016/j.renene.2018.09.005}}.

7. M. Abuella and B. Chowdhury,
"Improving Combined Solar Power Forecasts Using Estimated Ramp Rates: Data-Driven Post-Processing Approach,"
\textit{IET Renewable Power Generation}, vol.~12, no.~10, pp.~1127--1135, 2018.
DOI: \href{https://doi.org/10.1049/iet-rpg.2017.0447}{\textcolor{blue}{10.1049/iet-rpg.2017.0447}}.

\medskip
\textbf{Manuscript Under Review}

M. Abuella, A. Allahham, S. L. Walker,
"A Policy-Navigation Framework for Exploring Hydrogen Integration Pathways in Great Britain Towards Net Zero,"
Submitted to \textit{Energy Policy}, February 28, 2026.
Preprint available on SSRN:
\href{http://ssrn.com/abstract=6678965}{\textcolor{blue}{http://ssrn.com/abstract=6678965}}.

%--------------------------
\subsection*{Conference Papers}

1. M. Abuella, A. Allahham, S. L. Walker,
"Socio-Technical, Electrification, and Hydrogen-Driven Pathways for Residential Heating Decarbonisation in the North of Tyne,"
in \textit{9th International Conference on Environment Friendly Energies and Applications (EFEA 2025)}, 4--5 December 2025, Newcastle upon Tyne, UK.

2. M. Abuella, A. Allahham, S. L. Walker,
"A Game-Theoretic Framework for Hydrogen Integration to Accelerate Energy Transitions in Great Britain,"
in \textit{14th International Conference on Renewable Power Generation (RPG 2025)}, 24--26 October 2025, Shanghai, China.

3. M. Abuella, M. Atoui, S. Nowaczyk, S. Johansson, E. Faghani,
"Data-Driven Explainable Artificial Intelligence for Energy Efficiency in Short-Sea Shipping,"
in \textit{ECML PKDD}, 2023.

4. M. Abuella and B. Chowdhury,
"Adjusting Post-Processing Approach for Very Short-Term Solar PV Power Forecasts,"
in \textit{2021 IEEE MI-STA}, 2021.

5. M. Abuella and B. Chowdhury,
"Qualifying Combined Solar Power Forecasts in Ramp Events' Perspective,"
in \textit{IEEE Power and Energy Society General Meeting}, 2018.

6. M. Abuella and B. Chowdhury,
"Forecasting Solar Power Ramp Events Using Machine Learning Classification Techniques,"
in \textit{IEEE PEDG}, 2018.

7. M. Abuella and B. Chowdhury,
"Hourly Probabilistic Forecasting of Solar Power,"
in \textit{North American Power Symposium (NAPS)}, 2017.

8. M. Abuella and B. Chowdhury,
"Random Forest Ensemble of Support Vector Regression Models for Solar Power Forecasting,"
in \textit{IEEE ISGT}, 2017.

9. M. Abuella and B. Chowdhury,
"Solar Power Forecasting Using Support Vector Regression,"
in \textit{ASEM International Annual Conference}, 2016.

10. M. Abuella and B. Chowdhury,
"Solar Power Forecasting Using Artificial Neural Networks,"
in \textit{North American Power Symposium (NAPS)}, 2015.

11. M. Abuella and B. Chowdhury,
"Solar Power Probabilistic Forecasting by Using Multiple Linear Regression Analysis,"
in \textit{IEEE SoutheastCon}, 2015.

\end{document}
